\begin{table*}[t!]
\centering
\caption{Dataset properties and populations. Published and analyzed reference/sample counts are distinguished; the MOS column is the observed source-score range, retained without calibration or reversal. Modelling datasets influenced selection; evaluation datasets did not. Family identifies the four-dataset PCA-family study, whose populations equal the main analyzed populations; PointPCA is evaluated only on APSIPA because feature extraction for the other datasets exceeded the available 128 GB RAM. LS-PCQA uses 930 supplied accurate-MOS samples rather than its larger generated-score corpus.}
\label{tab_datasets}
\scriptsize
\begin{adjustbox}{max width=\textwidth}
\begin{tabular}{llllrrrrrll}
\toprule
ID & Dataset & Attributes & Distortions & Pub. refs & Pub. samples & Used refs & Used samples & MOS range & Role & Family \\
\midrule
1 & APSIPA \cite{apsipa_dataset} & Color & G-PCC, V-PCC & 8 & 254 & 8 & 232 & 1.00--5.00 & Modelling & Yes \\
2 & ICIP2023 \cite{basics_dataset} & Color, normals & Compression & 75 & 1494 & 45 & 898 & 1.12--4.82 & Evaluation & Yes \\
3 & LS-PCQA \cite{liu2023lspcqa} & Color & Noise, resampling, compression & 104 & $>22{,}000$ & 85 & 930 & 1.00--4.96 & Modelling & -- \\
4 & SJTU-PCQA \cite{yang2020predicting} & Color & Octree, resampling, noise & 9 & 378 & 9 & 378 & 1.00--9.80 & Evaluation & Yes \\
5 & UnB\_PC \cite{Torlig2018anovel} & Color, normals & Octree, JPEG & 6 & 60 & 6 & 60 & 1.20--5.00 & Evaluation & Yes \\
6 & WPC \cite{wpc} & Color & Noise, resampling, G/V-PCC & 20 & 740 & 20 & 740 & 1.81--94.30 & Modelling & -- \\
7 & WPC2 \cite{liu2021wpc2} & Color & V-PCC & 16 & 400 & 16 & 400 & 10.82--93.14 & Evaluation & -- \\
\bottomrule
\end{tabular}
\end{adjustbox}
\end{table*}
